\chapter{实践卷：并发、运行时和系统合同}

第二册实践卷分两层。第一层是本册已有实验：

\begin{lstlisting}[language=bash]
books/cpu-volume-2/labs/compute_systems
\end{lstlisting}

第二层是配套代码实践卷：

\begin{lstlisting}[language=bash]
books/compute-systems-engine-code/labs/reference_pipeline
\end{lstlisting}

本附录先把第一层并发实验拆成阶段。并发代码不能只看最终数字；它必须证明没有丢任务、重复执行、等待泄漏、关闭遗漏或把失败样本误当正常吞吐。

\section{零、贯穿主线工程}

第二册主线工程是 \filepath{compute_systems} 运行时实验，再向 \filepath{reference_pipeline} 代码实践卷扩展。第一层先把每章模型、单机并发和运行时合同做实：逐章 chapter models、串行 reference、并行归约、计数器、有界队列、wait channel、TaskRuntime、futex probe。第二层再把输入身份、manifest、split、merge 和恢复报告接上。不要把这些看成散装 demo；它们共同回答一个问题：一个计算系统怎样证明任务没有丢、没有重、没有卡住，并且能发布可复查报告。

\begin{longtable}{p{0.20\linewidth}p{0.30\linewidth}p{0.40\linewidth}}
阶段 & 主线代码 & 学到的工程能力 \\
\hline
逐章模型 & chapter models & 每章抽象概念都先落到可运行报告结构 \\
Reference 与切片 & 并行归约 & 串行结果是并行优化的可信尺子 \\
同步基础 & 同步模块 & mutex、atomic、队列和 close 都有可观察合同 \\
等待通道 & wait channel & 把唤醒事件和数据 drain 分开验证 \\
任务运行时 & TaskRuntime & submitted、completed、failed、rejected 和 CSV 报告可复查 \\
内核等待边界 & futex probe & 用户态快路径、内核慢路径和错误返回都要被测试 \\
恢复扩展 & 代码实践卷 & 进入输入身份、manifest、split、merge 和恢复报告 \\
\end{longtable}

\section{一、工程入口}

本册实验的主要文件如下：

\begin{lstlisting}[language=bash]
include/compsys/parallel_reduce.hpp
include/compsys/chapter_models.hpp
include/compsys/sync_primitives.hpp
include/compsys/task_runtime.hpp
include/compsys/wait_channels.hpp
include/compsys/futex_lab.hpp
tests/compsys_tests.cpp
src/chapter_models_probe.cpp
src/wait_channels_probe.cpp
src/task_runtime_probe.cpp
src/futex_lab_probe.cpp
\end{lstlisting}

构建和测试命令如下：

\begin{lstlisting}[language=bash,caption={第二册 compute systems 实验构建和测试}]
cmake -S books/cpu-volume-2 -B books/cpu-volume-2/build/compsys-debug -DCMAKE_BUILD_TYPE=Debug
cmake --build books/cpu-volume-2/build/compsys-debug
ctest --test-dir books/cpu-volume-2/build/compsys-debug --output-on-failure
\end{lstlisting}

\code{compsys\_chapter\_models\_probe} 是每章模型的人工入口。它不是替代测试，而是把 21 个章节/深入章节模型的报告字段按稳定顺序打印出来：

\begin{lstlisting}[language=bash,caption={运行第二册逐章模型探针}]
books/cpu-volume-2/build/compsys-debug/labs/compute_systems/compsys_chapter_models_probe
\end{lstlisting}

输出第一行是 \code{chapter\_models,21}。后续行分别对应第 1 章到第 18 章，以及第 5b、9b、12b 三个深入章。读者修改任一模型时，至少要同时更新三处：模型函数、\code{compsys\_tests.cpp} 中的断言、这个 probe 的输出字段。

\section{二、阶段 0：并行归约先对照串行 reference}

\code{sequential_sum} 是 reference，\code{parallel_sum} 是优化版。reference 是可信但不追求快的版本；并行版必须和它对同一份输入返回同一结果。

最小推导如下：把数组切成若干连续块，每个 worker 求一段局部和，最后主线程把局部和相加。若每个元素恰好属于一个块，且最终汇总不漏块、不重复块，那么并行结果等于串行结果。

失败样本也要写进报告：worker 数大于元素数时，不能生成空范围越界；worker 数为 1 时，应等价于串行；包含负数时，不能只测试自然数递增数组；空输入的语义必须明确。测试 \code{test_parallel_sum} 用含负数的固定数组对照 \code{sequential_sum} 和 \code{parallel_sum}。

\section{三、阶段 1：计数器和同步原语}

\code{MutexCounter} 和 \code{AtomicCounter} 都实现 \code{add/value}，但它们的保护方式不同。mutex 是互斥锁，一次只允许一个线程进入临界区；atomic 是原子变量，单个读改写操作不能被其他线程切开。

实践时不要只说“线程安全”。要写一个多线程样本：4 个线程，每个线程加 250 次，期望最终值是 1000。若少于 1000，说明更新丢失；若 value 读取没有同步，可能读到不一致状态。测试 \code{test_counters} 正是这条验收。

\section{四、阶段 2：有限队列和 backpressure}

\code{BoundedMpmcQueue} 是有容量上限的多生产者多消费者队列。MPMC 表示 multiple producers, multiple consumers。backpressure 是下游处理慢时，上游被容量限制住；它不是“队列满了”一句话，而是一个可观察合同：生产者等待次数、等待时间、关闭后返回值都要能报告。

队列实践按四步写：

\begin{enumerate}
  \item 单线程 \code{push/try_pop}：验证 FIFO 顺序、size 和 \code{max_size}。
  \item 消费者空队列阻塞后 \code{close}：验证 \code{pop} 返回空、\code{pop_closed} 增加、等待时间被记录。
  \item 生产者满队列阻塞后 \code{close}：验证第二次 \code{push} 返回 false、\code{push_closed} 增加。
  \item 多生产者多消费者 drain：容量为 1 时仍然不能丢值、重值或越界。
\end{enumerate}

这些测试对应 \code{test_queue}、\code{test_queue_close_wakes_consumer}、\code{test_queue_close_wakes_blocked_producer} 和 \code{test_queue_multi_producer_consumer_drain}。如果报告只写“push/pop 正常”，没有证明 close 唤醒，队列合同是不完整的。

\section{五、阶段 3：permit、barrier、future 和 atomic wait}

\code{run_permit_limiter} 训练“许可证”概念：最多允许 \code{permit_count} 个任务同时进入关键区域。测试要检查 \code{max_in_flight<=permit_count}，否则限流只是装饰。

\code{run_phase_barrier} 训练 barrier。barrier 是阶段边界：所有 worker 到达同一阶段后，才能一起进入下一阶段。报告中的 \code{arrival_count} 应等于 \code{worker_count*phase_count}，否则有人没到或被重复统计。

\code{run_future_result_paths} 训练 future 的三种结果：正常值、异常、broken promise。broken promise 指承诺方消失或没有设置结果，等待方不能永远挂住。

\code{run_atomic_wait_version_counter} 训练版本号等待。version 是单调增长的状态编号；等待者看到版本变化才继续。这里的 \code{wake_count} 和 \code{observed_sum} 证明等待路径真的被走到，而不是只靠忙等碰巧结束。

\section{六、阶段 4：wait channel 和 eventfd/epoll}

wait channel 是“等待者、唤醒条件和唤醒通道”的组合。只说“线程睡眠”不够，因为系统程序真正关心的是谁能唤醒、唤醒后读到什么、队列是否被完整 drain。

\code{run_eventfd_epoll_contract} 检查两类事实：eventfd 计数被读到几次，队列和信号量语义是否被 drain。报告字段 \code{epoll_ready_observations}、\code{epoll_empty_observations}、\code{queue_fifo_ok} 和 \code{semaphore_empty_after_reads} 都是合同的一部分。若只检查 epoll 返回“可读”，就还没证明数据结构本身正确。

\section{七、阶段 5：任务运行时和同池等待事故}

\code{TaskRuntime} 的报告字段包括 submitted、completed、failed、rejected、unfinished、max queue depth、active workers 和 accepting/stopping 状态。实践时要覆盖三类场景：正常完成、任务抛异常、shutdown 后提交被拒绝。

同池 future wait 是一个典型事故：worker 执行父任务，父任务又提交子任务并在同一个 worker 池中等待子任务；如果所有 worker 都被父任务占住，子任务没有 worker 可运行，就会饥饿。函数 \code{run_same_pool_future_wait_probe} 用报告字段描述这个停滞点，\code{run_continuation_runtime_probe} 则展示用 continuation 拆开等待边界后的路径。

这类实验的结论不能写成“future 不好”或“线程池不好”。正确结论是：同一执行池里阻塞等待自己提交的后续任务，会把 worker 资源耗尽；解决办法是改变调度结构、拆 continuation，或使用独立等待资源。

\section{八、阶段 6：futex 快路径和慢路径}

futex 的核心边界是用户态快路径和内核等待路径。无竞争时，锁应主要在用户态完成；有竞争时，等待者进入 futex wait，释放者 futex wake。报告中的 \code{contended_count}、\code{wait_count} 和 \code{wake_count} 用来证明慢路径真的发生过。

\code{run_futex_contract_probe} 还要覆盖三个容易被忽略的路径：值不匹配返回 \code{EAGAIN}，超时返回 timeout，被信号打断返回 \code{EINTR}。如果只测“最后计数正确”，就没有证明 futex 合同完整。

\section{九、和代码实践卷衔接}

当实验开始需要输入身份、诊断、manifest、checkpoint 或恢复报告时，应转入 \filepath{books/compute-systems-engine-code/labs/reference_pipeline}。本册原理解释机制，代码实践卷负责把机制做成可复查工程。连接规则很简单：原理卷中的对象如果已经有 C++ 类型，就引用类型名；如果还没有类型，就不能把它当作已实现能力。

\begin{longtable}{p{0.16\linewidth}p{0.28\linewidth}p{0.46\linewidth}}
阶段 & 代码切片 & 必须回答的问题 \\
\hline
0 & \code{sequential_sum/parallel_sum} & 每个元素是否恰好参与一次归约 \\
1 & \code{MutexCounter/AtomicCounter} & 多线程更新是否丢失 \\
2 & \code{BoundedMpmcQueue} & 满队列、空队列和 close 是否都能唤醒 \\
3 & permit 等同步对象 & 失败路径是否都被测试 \\
4 & eventfd 与 epoll & 唤醒后数据是否按合同 drain \\
5 & \code{TaskRuntime} probes & 是否有任务丢失、失败、拒绝或饥饿 \\
6 & \code{TaskRuntime} CSV & 报告字段是否能稳定导出和复查 \\
7 & futex probes & 快路径、慢路径和错误返回是否都可观察 \\
\end{longtable}

\begin{labbox}
大作业阶段 6 已经有入口：\code{write_task_runtime_report_csv} 能把 \code{TaskRuntime::Report} 导出为稳定 CSV。报告不能只测 happy path；至少要用正常完成、任务抛异常、关停后提交被拒绝、同池等待饥饿四类情况解释字段如何变化。
\end{labbox}

\section{十、每章代码任务矩阵}

第二册章节多、跨度大，但代码主线必须连续：先证明单机计算和同步合同，再进入运行时、异步 I/O 和分布式边界。每章读完都要把一个字段、一个状态或一个失败路径写成代码。

\begin{longtable}{p{0.18\linewidth}p{0.34\linewidth}p{0.38\linewidth}}
章节 & 已给出的完整模型代码 & 验收证据 \\
\hline
计算系统地图 & \code{model\_ch01\_input\_pipeline} & 输入总数、accepted、rejected 和 checksum 都由记录值推导出来 \\
核心、流水线、分支 & \code{model\_ch02\_branch\_predictor} & 交替输入下 actual taken 和 misprediction 不再只停留在概念 \\
Cache、TLB、虚拟内存 & \code{model\_ch03\_cache\_tlb} & 同一组 byte offset 同时推出 cache line、page 和 TLB miss \\
NUMA 与一致性 & \code{model\_ch04\_numa\_coherence} & local、remote 和 coherence transfer 分开计数 \\
测量、Roofline、计数器 & \code{model\_ch05\_roofline} & arithmetic intensity、memory roof 和 attainable GFLOPS 可断言 \\
PMU 与 SIMD 证据 & \code{model\_ch05b\_pmu\_simd\_evidence} & counter 可用性、retired instruction 和 vector fraction 同时进入报告 \\
布局、局部性、prefetch & \code{model\_ch06\_layout\_locality} & AoS/SoA 触碰字节数和 saved bytes 由字段数推导 \\
SIMD 与 ILP & \code{model\_ch07\_simd\_ilp} & vector iterations、scalar tail 和 lane utilization 可复查 \\
计算 kernel 模式 & \code{model\_ch08\_map\_filter\_reduce} & filter 后的 mapped sum 与 reduced sum 一致 \\
线程、调度、亲和性 & \code{model\_ch09\_thread\_scheduling} & worker 负载分配能暴露 imbalance \\
Linux 调度与内存管理 & \code{model\_ch09b\_linux\_wait\_wake} & wait loop 区分 spurious wakeup 和真正完成 \\
锁、条件变量、信号量 & \code{model\_ch10\_sync\_backpressure} & 生产、消费、容量和 backpressure event 都可观察 \\
原子与内存序 & \code{model\_ch11\_atomic\_publication} & stale read 和 monotonic observation 分开 \\
无锁结构 & \code{model\_ch12\_spsc\_ring} & push、pop、full rejection 和 empty pop 都有字段 \\
无锁证明与回收 & \code{model\_ch12b\_epoch\_reclamation} & protected node 延迟回收，未保护 node 才释放 \\
reduce、scan、partition & \code{model\_ch13\_reduce\_scan\_partition} & reduce sum、last prefix 和 partition true count 来自同一输入 \\
任务运行时 & \code{model\_ch14\_work\_stealing} & 初始队列不均衡时 stolen task 和 completed task 可见 \\
异步 I/O 与 backpressure & \code{model\_ch15\_async\_io\_backpressure} & submitted、completed、max in-flight 和背压事件同时验收 \\
分布式模型 & \code{model\_ch16\_distributed\_attempts} & logical task、duplicate attempt 和 stale generation 分开 \\
分区、复制、共识 & \code{model\_ch17\_quorum\_commit} & quorum size 和 committed 状态由 ack 集合推出 \\
分布式工程边界 & \code{model\_ch18\_checkpoint\_recovery} & checkpoint 数、replay 起点和 replay 数量明确 \\
\end{longtable}

这些函数都在 \filepath{include/compsys/chapter_models.hpp} 声明，在 \filepath{src/chapter_models.cpp} 实现，在 \filepath{tests/compsys_tests.cpp} 断言。完整源码见下一附录。这里的关键不是函数规模多大，而是每章至少有一个从输入到报告的闭合模型：输入是什么，状态怎样变化，哪些字段证明结论，错误输入怎样拒绝，都必须能被 C++ 代码说明。

\section{十一、递进大作业：Compute Systems Runtime}

\begin{longtable}{p{0.16\linewidth}p{0.30\linewidth}p{0.44\linewidth}}
阶段 & 交付物 & 通过标准 \\
\hline
0 & 逐章模型 & 21 个 \code{model\_ch*} 函数都能构建、测试并由 probe 打印 \\
1 & 串行 reference & 并行优化必须先和它对照 \\
2 & 并行归约 & 数据切片不丢、不重、不越界 \\
3 & 同步基础 & mutex、atomic、condition、semaphore 各有失败路径 \\
4 & 有界队列 & backpressure、close、drain 和报告字段可复查 \\
5 & wait channel & eventfd/epoll 的唤醒和数据 drain 分开验证 \\
6 & TaskRuntime & 正常、异常、拒绝、饥饿都能改变报告字段 \\
7 & CSV 发布 & \code{write_task_runtime_report_csv} 固定表头和字段顺序 \\
8 & 分布式衔接 & 需要 manifest 或 checkpoint 时接入代码实践卷 \\
\end{longtable}
